\subsection{Every copy uses a diagonal leaf cell}
\label{subsec:copy-count}

The body restrictions from Corollary~\ref{cor:body-roles} leave only one
possible source of copies: the two signs share the same positions at the
leaves.  We now locate that source at a fixed entry of the pattern.

\begin{proposition}\label{prop:copy-count}
For every integer $h\geq 1$, every ordered copy of $H$ in $A_h$ maps the
pattern entry $(s+1,s+1)$ into
\[
 \mathcal L_h
 =\bigcup_{p=0}^{m-1} R_h^+(p)\times C_h^+(p).
\]
In particular, $|\mathcal L_h|=m$ and
\begin{equation}\label{eq:copy-bound}
 N_H(A_h)\leq m n^{2k-2},
 \qquad
 \frac{N_H(A_h)}{n^{2k}}
 \leq \frac{2^{-h}}{d_h^2}.
\end{equation}
\end{proposition}

\begin{proof}
By Corollary~\ref{cor:body-roles}, the last two rows and columns of an
$H$-copy lie in the prescribed body sets of a single mode.  Their entries
must form
\[
 P=\begin{pmatrix}1&0\\1&1\end{pmatrix}.
\]
We check the four kinds of modes, using the order equivalences
\eqref{eq:tree-orders} and the entry rules \eqref{eq:variable-entries}.
Write $\pi(p)=\lfloor p/2\rfloor$ for
the parent index of a node.

In mode $V_i^+$, let $p$ be the child-level row index, $a$ the parent-level
row index, and $b$ the parent-level column index.  The order of the two
rows gives $\pi(p)\leq a$.  The bottom-right entry of $P$ gives $a\leq b$.
Consequently $\pi(p)\leq b$, so the top-right entry is $1$, a
contradiction.

In mode $V_i^-$, use the same names for the node indices at the same
levels.  The row order gives $\pi(p)\geq a$, while the top-right entry
$0$ gives $a\geq b$.  Thus $\pi(p)\geq b$, and the bottom-right entry is
$0$, again a contradiction.

In mode $W_{i,d}^+$, let $p,q$ be the child-level row and column indices,
and let $b$ be the parent-level column index.  The column order gives
$b\leq\pi(q)$, and the top-left entry $1$ gives $\pi(p)\leq b$.
Since $p\equiv q\equiv d\pmod 2$, it follows that $p\leq q$.  The
top-right entry is therefore $1$, contrary to $P$.

Finally, in mode $W_{i,d}^-$, the column order gives $\pi(q)\leq b$,
and the top-right entry $0$ gives $\pi(p)\geq b$.  The common parity
implies $p\geq q$.  If $i<h$, the top-left entry is
$\mathbf 1[p<q]=0$, which is impossible.  If $i=h$, that entry is instead
$\mathbf 1[p\leq q]$, because the leaf blocks are shared.  It can be $1$
only when $p=q$.  This places the top-left entry of the body, namely
$(s+1,s+1)$ in $H$, in $\mathcal L_h$.

Each leaf block has size $m/2^h=1$, so $\mathcal L_h$ contains exactly
$m$ cells.  For each such cell, fix the images of row $s+1$ and column
$s+1$ of $H$.  There are at most $n^{k-1}$ choices for the other row
indices and at most $n^{k-1}$ choices for the other column indices;
requiring increasing orders and the prescribed entries can only reduce
these numbers.  This proves the first inequality in
\eqref{eq:copy-bound}.  Since $n=d_hm$ and $m=2^h$, division by $n^{2k}$
gives the second.
\end{proof}


\begin{remark}\label{rem:hitting-not-repair}
Changing the $m$ entries of $\mathcal L_h$ from $1$ to $0$ destroys every
original copy, but creates new ones. Indeed, fix a leaf index $p$ and put
$d=p\bmod 2$. In mode $W_{h,d}^+$, choose the leaf row in $R_h^+(p)$
and any dummy row. Choose, in this order, a column in
$C_{h-1}^+(\pi(p))$ and the leaf column in $C_h^+(p)$.
These rows and columns have the required order. Their body is now
\[
 \begin{pmatrix}1&0\\1&1\end{pmatrix}:
\]
the upper-left entry is the unchanged parent--child entry, the
upper-right entry is the changed leaf diagonal, and the lower entries
are dummy incidences. Choosing one position from every anchor group
of this mode completes a new copy of $H$, since all anchor entries and
signatures are unchanged. Thus an original-copy hitting set need not
provide an $H$-free repair.
\end{remark}
